%==============================================================================
%  sections_v3/model_section.tex -- the model section of draft_v3.
%
%  SOURCE OF TRUTH: research/model_v4/MODEL_CARD.md, version stamp 2026-08-28
%  (re-review audit repairs, commit 59c0dfc).
%  Every statement below transcribes a
%  card row; where a proof file and the card differ, the card wins.  Card IDs
%  survive into the bracket titles so the card record, HANDOFF_sign.md and the
%  assembler can find each object by its ID.
%
%  Proofs live in sections_v3/proofs_section.tex; no proof environment appears
%  in this file.
%==============================================================================

\section{The two-round model}
\label{sec:model-v4}

\subsection{Object and partition}
\label{sec:object}

The disclosure rule \emph{is} the market's partition. A stake threshold $\tau$ and a filing window
$T$ split the histories that reach the control node into two cells: a \emph{flagged} cell, on which
the filing has landed before the control decision, and a \emph{pooled} cell, on which it has not.
Nothing else in the model performs that split, and the two cells are exclusive and exhaustive by
construction. The paper's question is how noise-trading intensity $\kappa$ moves bidder entry and
the expected takeover premium once the rule has drawn that line, and how the answer changes when the
line is redrawn.

The control-outcome object is the expected engagement-related premium
$\Dact(\kappa,\tau,T)$. Beside it sit three price-path objects, all of them produced by the model
rather than assumed: the run-up path $R_d$, the cumulative run-up $R$, and the filing-day jump $J$.
Lower $\tau$ is a tighter threshold margin; lower $T$ is a tighter window margin. The two margins are
separate policy coordinates and the results below treat them separately, because the model does not
deliver the same answer at both.

Three primitives are inherited rather than built here. Order flow is the discrete, ternary-noise
structure of the Kyle tradition \citep{Kyle1985}, in the tractable discrete form used by
\citet{EdmansGoldsteinJiang2015}; prices are competitive in the sense of
\citet{GlostenMilgrom1985}, set equal to the expectation of the terminal payoff given the public
history; and the takeover premium is the welfare object for dispersed minority shareholders for the
free-rider reason of \citet{GrossmanHart1980}. What is new is not any of these pieces but the
partition that sits between them, and the sections that follow state it exactly.

\subsection{Timing: the two rounds}
\label{sec:timing}

Trading runs over business days $d = 0,\dots,H$, with $H$ finite. The sequence is: pooled round
$\to$ flag or no flag $\to$ flagged round if applicable $\to$ bidder decision.

\begin{enumerate}[label=(\arabic*),leftmargin=2.2em,itemsep=3pt]
\item Nature draws the standalone value $v$ and the blockholder's private signal $s$. The
  blockholder picks one complete contingent plan $j$ from a finite ordered menu $\mathcal J$.
\item \emph{Round 1 --- pooled trading.} The plan's stake path executes over $d = 0,\dots,H$. Market
  makers see pooled order flow and set $P_d^P = \Eop[Y \mid \mathcal H_d^P]$.
\item \emph{Disclosure node.} The flag lands if and only if $D = 1$: the plan engages, crosses
  $\tau$ at some date $c < \infty$, and $c + T \le H$. The filing reveals $F = (B^F, a = 1)$.
\item \emph{Round 2 --- flagged trading, then the bidder.} If $D = 1$ the blockholder submits the
  flagged order $Q^F$, the market sets the flagged price $P^F = P(F,Q^F)$, and then the bidder
  decides. If $D = 0$ there is no flagged round and the bidder acts on the pooled history.
\end{enumerate}

Two features of this timing are load-bearing, and each is stated as a hypothesis rather than left in
the prose, because named proof steps fail without them.

\begin{assumption}[\S2 no-feedback: no within-window re-optimisation]
\label{asm:nofeedback}
The blockholder does not re-optimise inside the filing window. There is no feedback from realised
order flow or from realised prices into the executed path: $B_j(s,d)$, $q_{jd}(s)$ and $Q_j^F$ are
functions of $(j,s,d)$ and of $(j,s,\tau,T)$ alone.
\end{assumption}

\begin{assumption}[\S2 flag termination: the flag terminates the pooled round]
\label{asm:flagterm}
Pooled trading stops when the filing lands. The flagged round follows the filing, and the bidder
acts after the flagged round.
\end{assumption}

Assumption~\ref{asm:nofeedback} is what makes the executed path a function of the plan and the
signal, and Steps 3 and 6 of the conditional-independence argument behind Lemma~\ref{lem:L2} fail
without it. Assumption~\ref{asm:flagterm} is what makes $Q^F_j(s) = b^*_j(s) - B^F_j(s)$ the
blockholder's whole residual position; read the other way, the flagged-round step of
Proposition~\ref{prop:P1} fails.

\subsection{Primitives}
\label{sec:primitives}

\begin{definition}[\S4.1 primitives: values, signals, noise and the disclosure policy]
\label{def:primitives}
The standalone value is $v \sim N(\mu_v,\sigma_v^2)$ and the blockholder's private signal is
$s = v + \varepsilon$ with $\varepsilon \sim N(0,\sigma_\varepsilon^2)$ independent of $v$, so that
$\Eop[v \mid s] = \mu_v + \beta(s-\mu_v)$ with the Gaussian projection
$\beta = \sigma_v^2/(\sigma_v^2+\sigma_\varepsilon^2) \in (0,1)$. The bidder draws a private synergy
shock $\xi \sim N(0,\sigma_\xi^2)$, independent of $(v,s)$; $\bar S$ is mean bidder synergy and
$K > 0$ the bidder's entry cost. Takeover premia without and with engagement are $m_0$ and $m_1$,
with premium wedge $\Delta_m = m_1 - m_0 > 0$, and $\Delta_V \ge 0$ is the non-takeover value
created by engagement. Noise trading is ternary: $z_d \in \{-\bar z, 0, +\bar z\}$ with $\bar z > 0$,
$\Prb(z_d = 0) = 1-\kappa$ and $\Prb(z_d = \pm\bar z) = \kappa/2$, where $\kappa \in [0,1]$ is
noise-trading intensity. The disclosure rule is the pair $(\tau,T)$ with $T \in \{1,\dots,H\}$; the
blockholder's initial and maximum stakes are $b_0$ and $\bar b$ with $0 \le b_0 \le \bar b$.
\end{definition}

\begin{definition}[\S4.2 plans and the legal clock]
\label{def:plans}
The plan menu $\mathcal J$ is finite and ordered from least to most aggressive, with
$|\mathcal J| = J < \infty$. Plan $j$ carries an engagement flag $a_j \in \{0,1\}$, equal to $1$ for
Voice and $0$ for Exit and Hold. Its cumulative pooled stake at day $d$ is $B_j(s,d) \in [0,\bar b]$
with $B_j(s,-1) = b_0$; the terminal target is $b_j^*(s) = B_j(s,H)$. The threshold-crossing date is
$c_j(s;\tau) = \inf\{d : B_j(s,d) \ge \tau\}$, equal to $+\infty$ if the path never reaches $\tau$,
and the legal filing date is $f_j = c_j + T$. The disclosure indicator is
$D_j(s;\tau,T) = \ind\{a_j = 1,\ c_j < \infty,\ f_j \le H\}$, so that $D = 1 \Rightarrow a = 1$. The
stake at filing is $B_j^F = B_j(s,f_j)$ and the flagged-round order is
$Q_j^F = b_j^*(s) - B_j^F$. A finite ordered coarsening $\Gamma$ maps the stake increment to the
pooled order mark, $q_{jd}(s) = \Gamma\bigl(B_j(s,d) - B_j(s,d-1)\bigr)$, and observed pooled order
flow is $X_d = q_{jd} + z_d$.
\end{definition}

\begin{definition}[\S4.3 information, prices and the control outcome]
\label{def:prices}
The pooled public history is $\mathcal H_d^P = (X_0,\dots,X_d;\ \text{flag landed by } d)$, which is
finite. The filing message is $F = (B^F, a{=}1)$, and the flagged tuple is $F$ augmented by the
flagged order, $\Sfl = (B^F,Q^F,a{=}1)$. The control-node information set is the \emph{public}
information at the control node: $\mathcal I_H = \mathcal H_H^P$ on the pooled cell $\{D=0\}$ and
$\mathcal I_H = \Sfl$ on the flagged cell $\{D=1\}$. The bidder's own $\xi$ is private. Write
$\mathcal C_F$ and $\mathcal C_P$ for the flagged and pooled cells; they are exclusive and
exhaustive by construction. The engagement posterior is $\pi(\mathcal I) = \Prb(a=1\mid\mathcal I)
\in [0,1]$, equal to $1$ on $\mathcal C_F$, and the bidder-entry probability is
\begin{equation}
\label{eq:entry}
p(\mathcal I) \;=\; 1 - \Phi\!\left(\frac{P + K + m_0 + \pi\Delta_m - \bar S}{\sigma_\xi}\right)
\;\in\;(0,1).
\end{equation}
With $\mathsf B$ the entry indicator, the terminal shareholder payoff is
\begin{equation}
\label{eq:Y}
Y \;=\; (1-\mathsf B)\,(v + a\Delta_V) \;+\; \mathsf B\,(P + m_0 + a\Delta_m),
\end{equation}
and prices satisfy the inner fixed point $P(\mathcal I) = \Eop[Y\mid\mathcal I]$. Two conventions are
part of the model rather than notation. First, $P_{-1}^P := \Eop[Y]$ is the pre-trading pooled
price, which is needed whenever $c = 0$ --- as $T = H$ forces on every flagged history. Second,
$P_{\ND}(\mathcal H_{f^-}^P) := P_{f^-}^P$: the not-yet-disclosed price is the last pre-filing pooled
price at the \emph{same} realised order flow, whose history already carries ``flag not landed by
$f-1$'', and not a never-disclosed counterfactual. The price-path objects are the run-up path
$R_d = P_d^P - P_{c^-}^P$, the cumulative run-up $R = P_{f^-}^P - P_{c^-}^P$, and the filing-day
jump $J = P^F - P_{\ND}$; $R_d$, $R$ and $J$ are unsigned, and $J$ is not claimed to be
$\kappa$-invariant.
\end{definition}

The genuine fixed point sits at the control nodes. At an earlier pooled date $d < H$ the price is a
tower expectation of already-solved control-node values, with no self-reference; only the
control-node map is a fixed point to be solved.

\begin{definition}[\S4.3 objective: the blockholder's payoff $U_j$]
\label{def:U}
The blockholder's objective is the expected terminal value of the position the plan builds, net of
what it costs to build and to engage:
\begin{equation}
\label{eq:Uj}
U_j(s) \;=\; \Eop\Bigl[\, b_j^*(s)\,Y \;-\; \mathcal C_j^{\mathrm{trade}} \;-\; a_j\,C_j(s)
\;\Bigm\vert\; s,\, j \,\Bigr],
\end{equation}
where $\mathcal C_j^{\mathrm{trade}}$ is the execution outlay --- increments valued at the pooled
prices $P_d^P$ up to the plan's last pooled date, plus $Q_j^F(s)\,P^F$ when $D_j = 1$ --- and
$C_j(s) \ge 0$ is the engagement cost. Only two properties of \eqref{eq:Uj} are ever used:
\emph{plan-locality}, that $U_j$ depends on $j$ only through the executed stake path, the prices paid
on it, the terminal stake, the engagement flag and the cost; and \emph{integrability},
$\Eop[\max_{j}|U_j|] < \infty$ under Assumption~\ref{asm:A2p}.
\end{definition}

\begin{remark}[Timing convention for the engagement cost]
\label{rem:Ctiming}
Display \eqref{eq:Uj} does not date $C_j(s)$. The engagement cost may be booked either on completing
the plan or as sunk once the filing has landed. The two readings give the same round-2 comparison on
the round-2 deviation set, so the existence result of Proposition~\ref{prop:P1} does not depend on
the choice; the statement there is written so as not to commit to a reading.
\end{remark}

\begin{assumption}[TR: the \S4.1--\S4.3 primitive table restrictions]
\label{asm:TR}
The primitives of Definitions~\ref{def:primitives}--\ref{def:U} satisfy the following, which several
results below consume as a block.
\begin{enumerate}[label=(TR-\roman*),leftmargin=3.2em,itemsep=3pt]
\item \emph{Distributional forms and signs.} The distributions of Definition~\ref{def:primitives}
  hold as stated, with $\sigma_v^2,\sigma_\varepsilon^2,\sigma_\xi^2 > 0$, $K > 0$, $\bar z > 0$,
  $\Delta_m > 0$, $\Delta_V \ge 0$, and
  \begin{equation}
  \label{eq:m0}
  m_0 \;\ge\; 0,
  \end{equation}
  so that $\bar m(\mathcal I) = m_0 + \pi(\mathcal I)\Delta_m \ge 0$. Restriction \eqref{eq:m0} is
  what makes the inner pricing fixed point exist, be unique and be continuous; dropping it produces
  both nonexistence and three-root multiplicity in executed counterexamples. The initial stake
  satisfies $b_0 < \tau$: a pre-existing crossing is outside the core model.
\item \emph{Stake-path regularity.} For Voice plans $\partial_d B_j \ge 0$ and
  $\partial_s B_j \ge 0$; Hold paths are constant and Exit paths weakly decreasing in $d$. For every
  plan and every $d$, the map $s \mapsto B_j(s,d)$ is Borel. Borel regularity is automatic for Voice
  (monotone in $s$) and for Hold (constant) and is a genuine addition for Exit, where the primitives
  supply monotonicity in $d$ only; without it the pooled prices are not defined, because pooled
  pricing integrates over every type including Exit types. Stake levels are continuum-valued: the
  finiteness of Assumption~\ref{asm:A2p} covers the plan menu, the image of $\Gamma$, the noise
  support and the calendar, not the stake level.
\item \emph{Legal-clock objects.} The definitions of $c_j$, $f_j$, $B_j^F$, $Q_j^F$ and $b_j^*$ in
  Definition~\ref{def:plans} hold as stated, together with $D = 1 \Rightarrow a = 1$; $Q^F \ge 0$ for
  Voice plans; and, at fixed policies, $T' < T$ implies $B^F(T') \le B^F(T)$ and
  $Q^F(T') \ge Q^F(T)$.
\item \emph{Pricing and entry.} The terminal payoff is \eqref{eq:Y}, prices obey the convention
  $P(\mathcal I) = \Eop[Y\mid\mathcal I]$, entry obeys \eqref{eq:entry}, the control-node information
  set is the fill given in Definition~\ref{def:prices}, and the two price conventions
  $P_{-1}^P = \Eop[Y]$ and $P_{\ND} = P_{f^-}^P$ hold.
\end{enumerate}
\end{assumption}

\subsection{Equilibrium notion}
\label{sec:equilibrium-notion}

\begin{definition}[\S3: cutoff perfect Bayesian equilibrium]
\label{def:cpbe}
A \emph{cutoff perfect Bayesian equilibrium} is:
\begin{enumerate}[label=(\roman*),leftmargin=2.4em,itemsep=3pt]
\item a weakly ordered cutoff vector $k = (k_1 \le \dots \le k_{J-1})$ mapping the signal $s$ into a
  plan;
\item sequentially optimal pooled and flagged components;
\item Bayes-consistent beliefs on path;
\item competitive pooled and flagged prices at their fixed points;
\item the bidder-entry rule \eqref{eq:entry};
\item off-path beliefs as limits of full-support perturbations.
\end{enumerate}
The inequalities in (i) are weak, so collapsed action regions --- including a collapsed Hold region
--- are permitted. Existence is a Brouwer argument on the compact ordered polytope $\Theta$ for the
outer map $\Tmap(k;\vartheta)$ of Section~\ref{sec:polytope}, at a fixed point
$k = \Tmap(k;\vartheta)$. Uniqueness is \emph{not} claimed.
\end{definition}

\subsection{Standing hypotheses}
\label{sec:hypotheses}

The hypotheses below are the model's standing assumptions. They are not all maintained at once by
every result: each result in Section~\ref{sec:results} names the subset it consumes, and three of
them carry dated numerical records showing that they fail at the calibration the accompanying
implementation runs. Those records are stated with the assumptions they bear on, not deferred.

\begin{assumption}[A1: independent primitives]
\label{asm:A1}
$v$, $\varepsilon$, $\xi$ and all $z_d$ are mutually independent, and all variances are strictly
positive.
\end{assumption}

\begin{assumption}[A2$'$: finite model, amended boundedness]
\label{asm:A2p}
The plan menu $\mathcal J$, the image of $\Gamma$ (the order-mark support), the noise support
$\{-\bar z,0,+\bar z\}$ and the calendar horizon $H$ are finite. Prices and payoffs are locally
bounded in $(s,\vartheta)$ on the maintained parameter set, and
$\Eop\bigl[\max_{j\in\mathcal J}|U_j|\bigr] < \infty$ for every $k \in \Theta$.
\end{assumption}

The boundedness clause is stated as integrability because flat global boundedness is inconsistent
with the rest of the model: $v$ is Gaussian and the flagged region is unbounded in $s$ under
Assumption~\ref{asm:A7p}, so $Y$ --- and with it prices and $U_j$ --- is unbounded. Integrability is
all any proof below consumes.

\begin{assumption}[A3: ordered plans, single crossing]
\label{asm:A3}
At every belief and price system, adjacent-plan payoff differences cross zero at most once in $s$,
and the preferred plan is weakly increasing in $s$.
\end{assumption}

\begin{remark}[Numerical record for Assumption~\ref{asm:A3}]
\label{rem:A3record}
At the implemented calibration Assumption~\ref{asm:A3} itself fails, at two independently found
loci, upstream of Assumption~\ref{asm:A6}. First, at $(\kappa = 0.5,\ \tau_{50},\ T = 5)$ with $k_2$
on an \emph{open set} above cell edge $6$ --- verified at offsets $10^{-9}$ through
$2\times10^{-2}$ --- the difference $U_V - U_H$ has three strict sign changes, at
$s = 1.5754434$, $1.5833333$ and $1.5902426$, with middle excursions of
$2.4$--$2.8\times10^{-4}$ against a $10^{-9}$ payoff tolerance. The pointwise argmax runs
$H,V,H,V$, single-valued on each interval, so no weakly increasing selection exists: the selection
set is empty and the outer map $\Tmap$ is \emph{undefined} there, not merely discontinuous. Second,
at $(\kappa = 0.15,\ 0.05,\ 5)$ --- one of the four unresolved nodes of
Remark~\ref{rem:P1record} --- the argmax reverses from Voice to Hold across cell edge
$s = 1.659062163$ at both located fixed points, so the preferred plan decreases in $s$. The route is
the $s$-direction step of the Voice payoff, whose interior count $n(s)$ is integer-valued,
interacting with the off-path price snap. This does not conflict with the recorded finding that the
A3 and A6 proxies pass at every achieving seed: those proxies are local screens --- the A3 proxy
tests residual slope signs at the two candidate cutoffs and the A6 proxy tests non-pinning at the
closest seed --- and neither measures argmax monotonicity over $s$ nor continuity of $\Tmap$ in $k$,
so both are silent on these findings. A candidate mechanical account of the four unresolved nodes is
on file: at the $\kappa = 0.15$ node one fixed
point sits exactly on the edge where $U_H - U_V$ jumps through zero without crossing it, and the
panel's residuals, $3.06\times10^{-4}$ to $1.77\times10^{-3}$ at cutoff residuals of
$10^{-11}$ grade, bracket the recorded range exactly; the $k$-direction jump mechanism does
\emph{not} explain those nodes, and the panellist who proposed it recorded that negative himself.
That account has since been swept over the other three unresolved nodes under a pre-registered
three-way rule, and it holds at all three. At $(\kappa = 0.15,\ 0.075,\ 1)$ and
$(\kappa = 0.85,\ 0.075,\ 1)$ a located fixed point sits on an $n(s)$ cell edge --- at
$1.460178993$, offset about $10^{-13}$, where $10$ of the $30$ recheck seeds land, and at
$1.517932397$, offset about $10^{-12}$, reached by no seed and found only by the direct edge test
--- with $U_H - U_V$ again jumping through zero without crossing. Neither pin is its node's
achieving basin: their payoff residuals, $1.398\times10^{-3}$ and $1.314\times10^{-3}$, sit above
the recorded bests of $1.059\times10^{-3}$ and $3.061\times10^{-4}$, and each equals the larger
one-sided jump to at most $2.7\times10^{-4}$ relative. At $(\kappa = 0.85,\ 0.05,\ 5)$ no pin was
found at any candidate edge in $[1.29,\,2.11]$; the achieving basin's worst deviation instead sits
in the cell immediately above edge $1.583333333$, $0.0250\,\sigma_s$ from it, where the same jump
through zero occurs, at a deviation-to-jump ratio of $0.366$, inside the pre-registered factor of
three. Every pin is of the $n(s)$ family and the $\tau$-crossing pullbacks yielded none; the
distances of probe~5(b) replicate, at $0.0258/0.0437/0.0295\,\sigma_s$ against
$0.026/0.044/0.030$; and no node yields a second independent fixed point, so the residual bracket at
the first node does not recur. This is diagnostic evidence at one calibration, and existence at
these nodes stays neither claimed nor denied. No label moves on any of it: Assumption~\ref{asm:A3}
is a hypothesis, and Proposition~\ref{prop:P1} stays PROVED as a conditional. Records: the panel
files named in Remark~\ref{rem:A6record}, and
\texttt{quality\_reports/fixes/t2\_t34\_account\_sweep.py} with
\texttt{t2\_t34\_account\_sweep.json}.
\end{remark}

\begin{assumption}[A4: legal-clock discipline]
\label{asm:A4}
The crossing date $c$ is the first date the path reaches $\tau$; the filing lands exactly at $c+T$;
filings truthfully reveal stake and purpose; and only Voice plans cross in the core model.
\end{assumption}

\begin{assumption}[A5: inner pricing regularity, mostly demoted to a theorem]
\label{asm:A5}
Each public-history pricing map has a unique fixed point, continuous in beliefs, cutoffs and
parameters.
\end{assumption}

Assumption~\ref{asm:A5} is stated here for reference, and most of it is a theorem rather than an
assumption. Under \eqref{eq:m0}, existence, uniqueness and continuity of the \emph{inner} fixed point
are proved, not assumed. With the belief summaries $\bar y(\mathcal I) = \Eop_\mu[v] + \pi\Delta_V$
and $\bar m(\mathcal I) = m_0 + \pi\Delta_m$, the pricing map reduces to the scalar equation
$P \mapsto \bar y + \bar m\,p(P)/(1-p(P))$, which is strictly decreasing in $P$ wherever
$\bar m \ge 0$ and therefore crosses the identity exactly once. Three independent confirmations are
on file, one of them an executed counterexample producing zero roots and another producing three
once $m_0 < 0$. What Assumption~\ref{asm:A5} is retained for is its \emph{continuity} clause: the
pricing family is continuous in the cutoff vector and the parameters, and measurable in the flagged
tuple. Because the flagged information sets are continuum-indexed, ``unique fixed point'' must be
read as a measurably selected family, not a finite list. Where a result below cites
Assumption~\ref{asm:A5} for existence or uniqueness, it may cite \eqref{eq:m0} instead.

Two continuities must be kept apart here, since only one of them is proved. Continuity of the inner
root \emph{in its belief summaries} $(\hat v,\pi)$ follows from \eqref{eq:m0}. Continuity of the
\emph{composition} $k \mapsto (\hat v,\pi) \mapsto P$ in the cutoff vector is what the clause above
retains as an assumption: the $k$-dependence runs through the conditioning rather than through the
pricing map, and no step derives it. Remark~\ref{rem:A6record} records that composition measured
discontinuous at the implemented calibration, at the same loci and by the same checks, so this
clause carries adverse applicability evidence of its own. No label moves on that:
Assumption~\ref{asm:A5} is a hypothesis, and no result below consumes its cutoff clause.

\begin{assumption}[A6: compact outer self-map]
\label{asm:A6}
All best-response cutoffs lie in a common compact ordered polytope $\Theta$; the outer map $\Tmap$
of Section~\ref{sec:polytope} is continuous and maps $\Theta$ into itself.
\end{assumption}

\begin{remark}[Numerical record for Assumption~\ref{asm:A6}]
\label{rem:A6record}
% Long unbreakable repo paths below; scoped to this remark, no wording effect.
\setlength{\emergencystretch}{3em}%
The continuity clause of Assumption~\ref{asm:A6} fails for the declared construction of $\Theta$,
and the locus is not the one first proposed. All $k$-dependence of $U_j$ runs through the pooled
price vector, the flagged layer being $k$-free under Assumption~\ref{asm:A7J}; Bayes applies where
the reaching weight is positive but a $k$-free, plan-uniform posterior is imposed on the frontier,
so the price system can be discontinuous exactly on the boundary of the reaching set of each pooled
history. That set lies inside the union of the finitely many \emph{cell-edge hyperplanes}
$\{k_i = a\}$ and the \emph{collapse faces whose dying plan is the sole generator} of some reachable
pooled history --- not the collapsed cutoff vectors as such. The jump reaches $\Tmap$ with
non-vanishing weight, because $U_j$ integrates those prices against the deviator's own noise law,
with weight at least $\min(\kappa/2,\,1-\kappa)^{d+1}$, independent of the dying plan's population
mass; the vanishing-mass defusal is therefore refuted, by both panellists independently. The
largest-weakly-increasing-selection tie-break is pointwise in $k$ and passes the jump through, and no
$k$-independent perturbation family reconciles the limits: at fixed perturbation size the system is
continuous in $k$, and the discontinuity is created only in the limit, which the choice of family
cannot fix. On collapse faces proper, for menus with $J \ge 3$ in which a middle plan owns a
reachable exclusive pooled history entering some $U_j$, the interior limit
$\mu_v + \beta(c - \mu_v)$ varies over the face while any $k$-free family supplies one constant, so
continuity fails at every face point but at most one --- a continuum-face lemma derived in a single
panel pass and \emph{not} gate-checked. The implemented menu is not in that class: Exit and Hold pool
perfectly in order flow, and its Hold-collapse face is measured clean, with pooled prices within
$4.441\times10^{-16}$ as $k_1$ sweeps to full collapse. The payoffs $U$ are bit-identical across
that sweep; $\Tmap_2$ is not, moving $6.66\times10^{-16}$ --- three ulps --- at the one $k_1$ where
the price signature itself deviates most, so the invariance holds at the map's own root-finder
resolution rather than bit for bit. At the implemented
calibration the failure is live at the interior $n(s)$ cell edges instead: measured $\Tmap_2$ jumps
of $6.33\times10^{-3}$, $1.09\times10^{-2}$ and $2.83\times10^{-2}$ across steps in $k_2$ of at most
$2\times10^{-9}$ at $(\kappa = 0.5,\ \tau_{50},\ T = 5)$, measured independently by both panellists
with separate scripts and agreeing to three significant figures, with the belief snap matching the
predicted value to about $10^{-8}$ at all three edges at the truncation-and-cancellation crossover
bracket of $10^{-8}$ --- at the probes' own $10^{-9}$ bracket the first edge still holds, at
$4.0\times10^{-8}$, while the second and third are $1.2\times10^{-7}$ and $1.7\times10^{-7}$, which
is floating-point cancellation over a sliver $10^{-9}$ wide and not a gap in the prediction ---
surviving-type controls at about $3\times10^{-9}$, and
robustness at $1000$ times the breakpoint-merge tolerance; at $(\kappa = 0.15,\ 0.05,\ 5)$ the jumps
reach $0.16$ and a diagonal crossing of $\Tmap_2$ is destroyed. A chamber interior
$\Theta^+ = [1.23,\,1.245]\times[1.5253,\,1.5506]$ has been exhibited that is compact,
self-mapping and jump-free at the baseline node --- Brouwer runs on it verbatim and it contains
$k^\star$ --- but it is not the $\Theta$ the existence argument constructs, it cannot be exhibited
without approximately locating the fixed point first, and no such chamber exists at the
$\kappa = 0.15$ node, where a fixed point sits exactly on the edge $k_2 = 1.659062163$. No label
moves and none is licensed: A6 is a hypothesis, and Proposition~\ref{prop:P1} stays PROVED
as a conditional, in the same pattern as $\Atau$ --- what is on record is that its antecedent, read
with the $\Theta$ the proof constructs, is not satisfied by the implemented calibration. Two repairs
are on file, both outside the declared Brouwer-with-one-fixed-family route: a $t$-constrained game
with a Kakutani argument and $t \downarrow 0$, and a $k$-indexed concentration family, constructible
but with its $0/0$ corner unresolved. The implementation's off-path floor of $10^{-14}$ \emph{is} the
fixed-$t$ constrained game --- the standard repair already shipped, with the switch relocated by
about $10^{-9}$ rather than removed. Coverage: probes at one node per claim class plus a 27-node
census, not swept over $(\kappa,\tau,T)$. Nonexistence is neither claimed nor shown: $23$ of $27$
sweep nodes converge, and a discontinuous self-map may still have fixed points.

The three decisive measurements are now curated executed checks, each with a JSON verdict beside its
script. \texttt{t2\_a6\_edge\_jump\_check} replays both panellists' routes at their own filed
brackets and reproduces the three $\Tmap_2$ jumps, which agree across routes to a relative
$1.3\times10^{-4}$, with controls at $2.8$--$3.6\times10^{-9}$ and the $\pm10^{-6}$ robustness
intact. \texttt{t2\_a6\_node15\_check} reproduces the $0.1647$ jump, the destroyed crossing
($+1.0\times10^{-7} \to -6.70\times10^{-2}$) and the edge fixed point to $1.06\times10^{-12}$.
\texttt{t2\_a6\_collapse\_face\_check} reproduces the pooled prices within $4.441\times10^{-16}$.
Every figure these checks touch reproduces; two wordings of the panel record were corrected and the
numbers were not. The analytic weight bound $\min(\kappa/2,\,1-\kappa)^{d+1}$ is \emph{not} curated,
since no probe computes it; what is curated is its measured counterpart, the jump entering the
adjacent-plan payoff difference undiminished. Records:
\texttt{threads/\allowbreak 2026-08-27\_A6\_\allowbreak panel\_substantiate.md} and
\texttt{threads/\allowbreak 2026-08-27\_A6\_\allowbreak panel\_defuse.md}; the analysis-grade probes
behind them at
\texttt{quality\_reports/\allowbreak fixes/\allowbreak a6\_panel\_probes\_\allowbreak 2026-08-27/};
and the curated checks at
\texttt{quality\_reports/\allowbreak fixes/\allowbreak t2\_a6\_\allowbreak edge\_jump\_check},
\texttt{t2\_a6\_\allowbreak node15\_check} and \texttt{t2\_a6\_\allowbreak collapse\_face\_check},
each with its \texttt{.py} and \texttt{.json}.
\end{remark}

\begin{assumption}[A7: filing sufficiency]
\label{asm:A7}
On flagged histories, $(B^F,Q^F,a{=}1)$ identifies the informed component of the selected plan;
conditional on it, the pooled order-flow residual is pure noise, independent of $(v,s,\xi)$.
\end{assumption}

This weak identification wording is not enough for Lemma~\ref{lem:L2}: it permits two pairs $(j,s)$
with different pooled paths, which is exactly that lemma's first failure case. Two injective forms
are therefore named separately, and the form is named every time either is used.

\begin{assumption}[A7$'$: on-path composed target]
\label{asm:A7p}
At a fixed cutoff policy, the composed terminal target $s \mapsto b^*_{j(s)}(s)$ is strictly
increasing on the flagged signal region, and this holds for every cutoff vector $k \in \Theta$.
Strictness is required only for flag-capable composed targets: passive plans that never flag need not
have strictly increasing $b_j^*$, and there is no backtracking of $b_j^*$ across admissible
Voice-plan switches.
\end{assumption}

\begin{assumption}[A7-J: joint tuple injectivity]
\label{asm:A7J}
The full map $(j,s) \mapsto (B_j^F,\,Q_j^F,\,a_j)$ is injective on the flagged-pair set
$\{(j,s) : D_j = 1\}$, including flagged pairs that are not selected on path.
\end{assumption}

Assumption~\ref{asm:A7J} is strictly stronger than Assumption~\ref{asm:A7p}, and it is the form the
existence argument consumes. Strictness of $B^F$ alone is neither necessary --- it fails at
crossing-date jumps on the pro-rata menu --- nor sufficient, because of multi-Voice backtracking.
Under Assumption~\ref{asm:A7p} the flagged tuple is continuum-valued as a tuple, since injectivity
forces $(B^F,Q^F)$ to be continuum-valued, though the coordinates may trade the burden between them;
injectivity plus measurability already gives the measurable inverse on standard Borel spaces, so no
separate assumption is needed for that. Satisfiability is resolved for Assumption~\ref{asm:A7p}:
together with a fixed cutoff policy and $\Omega > 0$ it delivers the on-path injective form on
positive-probability flagged tuples with an explicit inverse, and a satisfying menu exists --- the
pro-rata single-Voice menu with terminal target strictly increasing on all of $\mathbb R$, which also
satisfies Assumption~\ref{asm:A7J}. Assumption~\ref{asm:A7J} additionally needs $b^*$ strictly
increasing off the Voice region: a target flat below the Voice cutoff breaks it, in an executed
check producing forty collisions, while leaving Assumption~\ref{asm:A7p} intact. The failure boundary
is a binding stake cap, quantized stakes, a composed target repeating values
across Voice-plan switches, $\Omega = 0$, and policy dependence when the condition is stated only at
one equilibrium's cutoffs. Menus satisfying Assumption~\ref{asm:A7p} are fully separating on the
flagged set, so the burden moves to the incentive compatibility of Proposition~\ref{prop:P1}, not
away.

\begin{assumption}[A8: interior crossing]
\label{asm:A8}
$0 < \Omega(\kappa,\tau,T) < 1$. This is required only for positive cell mass, never for the
structural partition.
\end{assumption}

\begin{assumption}[$\Atau$: threshold chord restriction]
\label{asm:Atau}
The pooled posterior law has the symmetric ternary representation
\begin{equation}
\label{eq:atau}
\Eop[h] \;=\; A_0(\kappa)\,h(0) \;+\; A_{1/2}(\kappa)\,h(\bar\pi/2) \;+\; A_1(\kappa)\,h(\bar\pi),
\end{equation}
with $A_0' = A_1' = \Akap$ and $A_{1/2}' = -2\Akap$, and maintained orientation
$C_h(\bar\pi) \le 0$ with $|C_h|$ weakly increasing in $\bar\pi$. The strict version of the
orientation is the frozen manuscript's condition (C\textasteriskcentered); the case $C_h = 0$ must be
handled explicitly. Two further clauses are part of the restriction:
\begin{enumerate}[label=($\tau$-\roman*),leftmargin=3.2em,itemsep=3pt]
\item \emph{The kernel depends on the information set only through the engagement posterior:}
  $h(\mathcal I) = h(\pi(\mathcal I))$, so the three numbers $h(0)$, $h(\bar\pi/2)$, $h(\bar\pi)$ are
  well defined and $\kappa$-free. This is a restriction, not a reading. In the model
  $h = \pi\,p(\hat v,\pi)$ is a function of \emph{two} scalars, because \eqref{eq:entry} makes entry
  depend on the price as well as on the posterior; the clause says the standalone-value channel and
  the engagement channel do not co-move inside the pooled cell in a way that moves $h$ at a fixed
  posterior.
\item \emph{The support and $\bar\pi$ are $\kappa$-free; only the weights move.} The three points
  $\{0,\bar\pi/2,\bar\pi\}$ do not vary with $\kappa$, and $\bar\pi$ itself is $\kappa$-free at fixed
  $(\tau,T)$. Without the second half the conclusion of Lemma~\ref{lem:L3} is false: the derivative
  gains a term that is first order in $\bar\pi$ and the vanishing fails.
\end{enumerate}
\end{assumption}

Where the bite of Assumption~\ref{asm:Atau} actually sits is worth stating, because it is narrower
than the display suggests. Given a $\kappa$-invariant three-point support, the derivative
restrictions $A_0' = A_1' = \Akap$ and $A_{1/2}' = -2\Akap$ are not an extra assumption at all: they
are \emph{equivalent} to $\kappa$-invariance of the pooled block's total mass and of its
unnormalised engagement moment, both of which the model delivers at fixed policies. The entire
remaining content of Assumption~\ref{asm:Atau} is the support condition. A one-round ternary-noise
market with informed mark $2\bar z$ and pre-order engagement share $\tfrac12$ satisfies it; the
frozen manuscript's own no-disclosure structure, with informed mark $\bar z$, does not, since its
pooled law has four atoms, two of which move with $\kappa$. Whether the two-round pooled cell of
Section~\ref{sec:timing} satisfies the support condition is \textbf{open}, and every result
conditional on Assumption~\ref{asm:Atau} --- Lemma~\ref{lem:L3}, leg 3 of Lemma~\ref{lem:L4}, and
Part~(B) of Theorem~\ref{thm:T1} --- inherits that conditionality.

\begin{remark}[Numerical record for Assumption~\ref{asm:Atau}]
\label{rem:AtauRecord}
At the implemented calibration Assumption~\ref{asm:Atau} \textbf{fails}. The decisive representation
failure is already established by the support condition alone; the derivative pattern also fails, and
independently. The support half carries the verdict because it is the entire remaining content of
the restriction, and of the derivative-pattern failure only the $A_{1/2}'$ residual is inherited
from the support. The pooled cell's engagement-posterior law was enumerated exactly --- all
$4^{H+1} = 4{,}194{,}304$ order-flow paths, the same law the implementation integrates --- at $200$
nodes: $\kappa \in \{0.05,\dots,0.95\}$ crossed with the five frozen $\tau$ percentiles and
$T \in \{1,2,5,10\}$, at frozen policies with $H = 10$. Two gates pass first, so the object measured
is the one Assumption~\ref{asm:Atau} is about: an independent re-enumeration reproduces the pooled
mass to $0.0$ exactly, and the enumerated mean equals the pooled share
$\bar\pi_{\mathrm{pr}} = \Prb(a=1\mid D=0)$ to $1.7\times10^{-16}$. Neither a particular value of
$|\Akap|$ nor level symmetry is imposed anywhere, and $\bar\pi$ is read as the upper support point
throughout. Twenty nodes are degenerate --- $\bar\pi_{\mathrm{pr}} = 0$ at $T \in \{1,2\}$ with
$\tau$ at the tenth percentile, where no engaging atom survives into the pooled cell, the law is the
point mass at $0$, $M_P = 0$ and $C_h(0) = 0$, so the restriction holds vacuously and the node
decides nothing. At all $180$ non-degenerate nodes $\Atau$ fails; at none does it hold.
\begin{itemize}[leftmargin=1.6em,itemsep=3pt]
\item \emph{Clause ($\tau$-ii), support half --- fails, by some eleven orders of magnitude.} The
  support carries $23$ to $767$ distinct posterior values, never three ($0$ of $180$ nodes), and
  there is no mass at $\bar\pi/2$ at any node, so $A_{1/2} \equiv 0$. Between $0.57\%$ and $91.8\%$
  of the pooled mass sits off $\{0,\bar\pi/2,\bar\pi\}$, and $13.9\%$ at the median node
  ($T = 5$, median $\tau$, $\kappa = 0.55$: $107$ atoms, $A_0 = 0.768$, $A_1 = 0.093$). The atoms are
  not dust: coarsening the cluster tolerance to $10^{-3}$ still leaves $6$ to $332$ of them, and the
  floor-free law counts at most $51$ atoms fewer than the floored law the package prices. The
  interior atoms move with $\kappa$: the two-sided Hausdorff distance between adjacent-$\kappa$
  support sets reaches $0.4608$, unchanged when restricted to atoms carrying mass at least
  $10^{-6}$, against a predicted $<10^{-12}$, at $0$ of $18$ series.
\item \emph{Clause ($\tau$-ii), $\bar\pi$ half --- holds.} $\bar\pi = 1$ to $1.5\times10^{-13}$ at
  every non-degenerate node, and $\kappa$-free to the same order, at $18$ of $18$ series. This is a
  separate finding and not a partial rescue: $\bar\pi = 1$ is the one-round outcome, and the
  conjecture that the two-round timing leaves the pooled cell with a top atom strictly below $1$ is
  false at this calibration, because unflagged Voice types still generate fully revealing order
  flows. Across the whole grid $\bar\pi \in \{0,1\}$ and never interior, so the small-$\bar\pi$
  corollary of Lemma~\ref{lem:L3} has no instance here either.
\item \emph{Derivative pattern --- fails, and independently of the support.} $A_0' = A_1'$ holds at
  $0$ of $180$ nodes: $|A_0' - A_1'| \in [0.041,\,2.306]$ against a predicted $<10^{-10}$, with
  $A_0' \in [-2.146,\,2.374]$ against $A_1' \in [-0.014,\,0.429]$, an order of magnitude apart in
  level, and both change sign over the grid, which independently corroborates that $\Akap$ carries
  no sign. The restriction $A_{1/2}' = -2\Akap$ also fails at all $180$, but with
  $A_{1/2} \equiv 0$ that residual is exactly $2|A_0'|$ and is recorded as inherited --- a
  restatement of the support failure, not a second piece of evidence.
\item \emph{Chord identity --- fails.} The residual
  $\bigl|\mathcal S_P - \Delta_m\,|\Akap|\,|C_h(\bar\pi)|\bigr|$, with $\Akap$ recovered from the
  enumerated weights and $\bar\pi$ the actual upper support point, is $0.0013$ to $0.0717$ --- up to
  $7.17$ premium percentage points --- against a predicted $<10^{-10}$, at $0$ of $180$ nodes and on
  the most favourable of three kernel conventions. Recovered $|\Akap| \in [0.042,\,2.374]$; the value
  the identity would require is $[0.00023,\,0.392]$, disjoint from the separately implied
  $[0.997,\,1.158]$, which is a different object (a mean absolute slope over the $\kappa$ grid, under
  level symmetry) --- and the distance between the two measures what the level-symmetry assumption
  was doing.
\item \emph{Clause ($\tau$-i), reported as a diagnostic and not part of the verdict.} Within a
  cluster on which the posterior is constant to $10^{-12}$, the enumerated entry probability still
  spreads by up to $0.085$ and $h$ by up to $0.018$ mass-weighted. The kernel does not reach the
  information set only through the posterior at this calibration either.
\end{itemize}
This is NUMERICAL-class \emph{applicability} evidence at one calibration. No label moves,
and none is licensed: $\Atau$ is an assumption, not a labelled claim. Lemma~\ref{lem:L3}, leg 3 of
Lemma~\ref{lem:L4} and Part~(B) of Theorem~\ref{thm:T1} stay PROVED as conditionals with
their proofs untouched; what is now on record is that their antecedent is not satisfied by the
implemented pooled cell at this calibration, so at this calibration those legs say nothing about the
implemented cell. The open question above is a question about the assumption's \emph{domain}: a
different menu, a different $H$ or a different calibration could still satisfy the support condition.
Coverage caveats carried forward: the $18$ non-degenerate series are only \emph{six distinct pooled
cells}, because $T = 1$ and $T = 2$ induce identical partitions at every $\tau$, $T = 5$ joins them
at the three highest $\tau$ percentiles and repeats itself at the two lowest, and all five $T = 10$
quantiles coincide --- and all six fail; and the fifty $T = 10$ nodes sit at $\Omega = 0.000681$,
below the implementation's minimum cell mass. Script and record:
\texttt{quality\_reports/fixes/t2\_atau\_support\_check.py} $\to$
\texttt{t2\_atau\_support\_check.json} ($200$ nodes, $920$ pooled enumerations, $1002$ seconds;
top-level verdict field \texttt{FAILS at calibration}).
\end{remark}

\begin{assumption}[$\Abr$: chord--sensitivity bridge]
\label{asm:Abr}
For two compared thresholds $\tau' < \tau$ at fixed policies and a common $\kappa$:
\begin{enumerate}[label=(br-\roman*),leftmargin=3.2em,itemsep=3pt]
\item \emph{Representation at both policies.} The symmetric ternary representation \eqref{eq:atau}
  holds for the pooled class under $\tau$ \emph{and} under $\tau'$, with chord endpoints
  $\bar\pi(\tau)$ and $\bar\pi(\tau')$ and weight-derivative coefficients $\Akap(\tau)$ and
  $\Akap(\tau')$.
\item \emph{$\kappa$-localisation.} At fixed policies all $\kappa$-dependence of $M_P$ sits in the
  weights of \eqref{eq:atau}: the three support points $\{0,\bar\pi/2,\bar\pi\}$ and the kernel $h$
  \emph{as a function of the posterior} do not move with $\kappa$. Hence
  $\partial_\kappa M_P = \Delta_m \Akap C_h(\bar\pi)$ exactly, with no
  composition-through-$\kappa$ remainder. Against the literal display \eqref{eq:atau} this would
  restate (br-i); it is written against the honest reading $h = \pi\,p(\hat v,\pi)$, and it repairs
  that ambiguity rather than adding a fourth independent restriction, naming the same object as
  clause ($\tau$-i). The trailing ``hence'' is derivable, not assumed.
\item \emph{Coefficient stability across the threshold margin.}
  $|\Akap(\tau')| \le |\Akap(\tau)|$. The weakest sufficient form is equality: reclassification
  changes \emph{which} histories are pooled, not the $\kappa$-responsiveness of the pooled weights.
\item \emph{Endpoint linkage.} $\bar\pi$ is the chord endpoint of \eqref{eq:atau} --- the upper
  support point of the pooled posterior law --- and it is a weakly increasing function of the pooled
  prior engagement share $\bar\pi_{\mathrm{pr}} = \Prb(a=1 \mid D=0)$, \emph{the same function at
  $\tau$ and at $\tau'$}. The identity branch $\bar\pi = \bar\pi_{\mathrm{pr}}$ is excluded as
  degenerate.
\item \emph{Comparability of the chord functional across thresholds.} The chord functional
  $C_h(\cdot)$ --- and the kernel $h$ it is built from --- are the same functions of the posterior at
  both compared thresholds. Without it, leg 3 of Lemma~\ref{lem:L4} compares $|C_h|$ across two
  different functionals and the comparison is meaningless: $h = \pi p$ with $p$ priced off a cell
  whose composition the threshold moves, so $\tau$-invariance of $h$ is real content, not
  bookkeeping.
\end{enumerate}
\end{assumption}

Assumption~\ref{asm:Abr} is consumed by leg 3 of Lemma~\ref{lem:L4} and by Part~(B) of
Theorem~\ref{thm:T1}, and by nothing else. Clause (br-v) was independently required by three agents,
one of whom confirmed that it is not implied by (br-i)--(br-iv). One sharpening is recorded rather
than assumed: with $\rho := \tfrac12 A_{1/2} + A_1$, which is provably $\kappa$-free,
$\bar\pi = \bar\pi_{\mathrm{pr}}/\rho$, so (br-iv) is equivalent to
$\rho(\tau')/\rho(\tau) \ge \bar\pi_{\mathrm{pr}}(\tau')/\bar\pi_{\mathrm{pr}}(\tau)$; under the
level-symmetric reading $\rho = \tfrac12$ and $\bar\pi = 2\bar\pi_{\mathrm{pr}}$, which forces
$\bar\pi_{\mathrm{pr}} \le 1/2$, an inherited restriction on the domain of
Assumption~\ref{asm:Atau} that Lemma~\ref{lem:L4} does not resolve. Clause (br-iii) is the clause
with the least justification behind it, and it is the one to attack first. Because (br-i) carries the
representation \eqref{eq:atau} at both thresholds, the record in Remark~\ref{rem:AtauRecord} bears
directly on Assumption~\ref{asm:Abr} as well.

\begin{assumption}[AGE: general-equilibrium differentiability and contraction]
\label{asm:AGE}
On a candidate region $\mathcal R$ the outer map is twice continuously differentiable,
$L_{\mathcal R} < 1$, and the sign of the equilibrium liquidity derivative is constant on
$\mathcal R$.
\end{assumption}

\subsection{The cutoff polytope and the outer map}
\label{sec:polytope}

\begin{definition}[\S4.5 equilibrium objects: $\Theta$, $\Tmap$, and the general-equilibrium bounds]
\label{def:theta}
The cutoff vector is $k = (k_1,\dots,k_{J-1})$ with $k_1 \le \dots \le k_{J-1}$; when the menu is the
four named actions of the frozen manuscript it maps to that draft's triple of cutoffs. The cutoff
polytope $\Theta$ is nonempty, compact and convex, and $\vartheta$ is the parameter vector. The
\emph{outer cutoff best-response map} is $\Tmap(k;\vartheta)$, a continuous self-map of $\Theta$; it
is always written calligraphically, since upright $T$ is the filing window. On a region $\mathcal R$
the contraction bound is $L_{\mathcal R} = \sup_{\mathcal R}\lVert D_k\Tmap\rVert$, required to be
below $1$ by Assumption~\ref{asm:AGE}. The strictness coordinates are $r_\tau = -\tau$ and
$r_T = -T$, so that higher $r$ is tighter. The direct fixed-policy attenuation margin is
\begin{equation}
\label{eq:gPE}
g_r^{PE} \;=\; -\operatorname{sgn}\!\bigl(d\Dact/d\kappa\bigr)\;\partial_{\kappa r}\Dact,
\end{equation}
with the sign written inline rather than carried by a symbol. The inversion-free derivative bounds
are $\bar k_x = |\partial_x\Tmap|/(1-L_{\mathcal R})$ and $\bar k_{\kappa r}$, and the
general-equilibrium remainder bound is
\begin{equation}
\label{eq:BGE}
\mathcal B_r^{GE} \;=\; |\Delta_{\kappa k}|\,\bar k_r
\;+\; \bigl(|\Delta_{kr}| + |\Delta_{kk}|\,\bar k_r\bigr)\bar k_\kappa
\;+\; |\Delta_k|\,\bar k_{\kappa r}.
\end{equation}
The \emph{dominance-and-contraction region} is $\mathcal R_r$, with slack
$\eta_r = g_r^{PE} - \mathcal B_r^{GE}$; the region may be empty.
\end{definition}

\begin{definition}[\S4.4 premium decomposition and comparative-statics objects]
\label{def:premium}
The engagement-premium kernel is $h(\mathcal I) = \pi(\mathcal I)\,p(\mathcal I)$, with $h \ge 0$ and
$h(0) = 0$, and the expected engagement-related premium is
\begin{equation}
\label{eq:Dact}
\Dact \;=\; \Delta_m\,\Eop\bigl[h(\mathcal I_H)\bigr] \;\ge\; 0.
\end{equation}
The cell averages are $M_F = \Delta_m\Eop[h \mid D=1]$ and $M_P = \Delta_m\Eop[h \mid D=0]$, each
defined when its cell has mass. The unconditional flagged weight is $\Omega = \Prb(D=1) \in [0,1]$,
which factors as $\Omega = \Prb(a=1)\,\omega_a$ with $\omega_a = \Prb(D=1 \mid a=1)$ the disclosed
share of engagements. The quantity $\bar\pi$ is the \emph{upper support point} of the pooled
engagement posterior in the representation \eqref{eq:atau}. The pooled engagement share is the
\emph{mean} of that law, not $\bar\pi$; under Assumption~\ref{asm:Atau} the share is
$\kappa$-invariant, being a mean-preserving spread, so it cannot be the quantity whose
$\kappa$-motion Lemma~\ref{lem:L3} describes, it is strictly below $\bar\pi$ in any non-degenerate
case, and it equals $\bar\pi/2$ only under the level symmetry $A_0 = A_1$. The liquidity
sensitivities are $\mathcal S = |\partial_\kappa\Dact|$ and $\mathcal S_P =
|\partial_\kappa M_P|$. The chord is
\begin{equation}
\label{eq:chord}
C_h(\bar\pi) \;=\; h(0) \;-\; 2h(\bar\pi/2) \;+\; h(\bar\pi),
\end{equation}
maintained non-positive with $|C_h|$ weakly increasing in $\bar\pi$, and $\Akap$ is the common
derivative of the weights in \eqref{eq:atau}, bounded on $[0,1]$. The weight-effect ratios are
$W_\tau$ and $W_T$ --- for instance $W_T = (1-\Omega(\tau,5))/(1-\Omega(\tau,10))$, at most $1$ when
$\Omega$ rises --- and the composition-effect ratios are $C_\tau$ and $C_T$, for instance
$C_T = \mathcal S_P(\tau,5)/\mathcal S_P(\tau,10)$, which are unsigned. The margin subscript on a
composition ratio is always written, because $C$ is otherwise overloaded by the chord $C_h$, the
engagement cost $C_j(s)$ and the cells $\mathcal C_F,\mathcal C_P$.
\end{definition}

Reading $\bar\pi$ as the mean of the pooled posterior law rather than as its upper support point
forces a point mass at $\bar\pi$ with $\Akap = 0$ and zero interior motion for every kernel. That is
degenerate, and it is excluded throughout.
